\documentclass[11pt]{article}
\usepackage[a4paper,margin=23mm]{geometry}
\usepackage{amsmath,amssymb,amsthm,mathtools,bm,booktabs,microtype,xurl,hyperref}
\hypersetup{colorlinks=true,linkcolor=blue,urlcolor=blue,citecolor=blue,
pdftitle={P54 Goldstone IR audit and background control}}
\setlength{\emergencystretch}{2em}
\newcommand{\tr}{\operatorname{tr}}
\newcommand{\Tr}{\operatorname{Tr}}
\newcommand{\diag}{\operatorname{diag}}
\newcommand{\MS}{\overline{\mathrm{MS}}}
\newtheorem{proposition}{Proposition}
\title{Route-F P54: Goldstone Infrared Audit,\\
Scalar Momentum Dependence and Background Control}
\author{P54PQ-v2 four-dimensional mainline}
\date{13 September 2026}
\begin{document}
\maketitle
\begin{abstract}
We retain the P54PQ-v2 action and first audit its frozen scalar point.
All 34 Goldstone hard mass shifts cancel against the same finite
fixed-VEV counterterms by the full orbit Ward identity. The excluded
soft scalar tadpoles vanish, but their radial potential Hessian has a
nonzero logarithmic infrared singularity. We compute its actual Gram
coefficient and the complete one-loop scalar momentum dependence on
the three-dimensional radial slice. The hard scalar kinetic correction
has a largest generalized eigenvalue 2.61, independently exposing a
failure of a small-insertion expansion. We then analytically rule out
an entire uniformly weakened bosonic parameter ray at fixed radii.
Four less-hierarchical scalar engineering points are tested separately;
none passes both radial positivity and perturbative-control criteria.
No physical benchmark, flavor fit, full gauged pole kernel, or complete
finite matching is claimed. The original inputs are not overwritten.
\end{abstract}

\section{Fixed action, conventions and acceptance boundary}
The scalar content is real $54$, complex $126$, complex $10$ and complex
$S$, with three chiral $16_F$ families. The original radial point is
\begin{equation}
 r_0=(w,\sigma,v_s)=(1,0.1265,0.25)\omega,
 \qquad \mu_0=g_U\omega=0.5626028899853225\omega.
\end{equation}
The same DR/$\MS$ background-field Landau convention and common
fixed-VEV mass counterterms of P-REN1 are used. The 328-dimensional
real canonical scalar Hessian has 38 soft tree modes: 33 gauge
Goldstones, one physical PQ mode, and four components of the tuned
light doublet. There are 290 positive hard scalar modes and 33 massive
gauge bosons. The 12 unbroken vectors remain massless along this radial
slice. Their vanishing potential contribution does not eliminate all
momentum-dependent gauge diagrams.

Let $R$ embed $r$ in the canonical field $X$, so
\begin{equation}
 G=R^TR=\diag(12/5,2,1),\quad
 \mathsf H=V_0''(X),\quad T_r=\partial_r\mathsf H,\quad
 U_{rs}=\partial_r\partial_s\mathsf H.
\end{equation}
All these are obtained from the actual polynomial action. Its quartic
degree makes $\mathsf H$ quadratic, hence cached radial jets determine
its exact radial dependence without a new Hessian evaluation or a fitted
mass formula. Floating-point residuals only test these identities.

Goldstone resummation and a finite physical scalar mass are distinct
questions. The former does not by itself render the second derivative
of a zero-momentum potential finite; external-momentum self-energies
are needed \cite{bg}. Hard/soft separation also matters in organizing
the resummation \cite{ek}. We apply these distinctions rather than
introduce a positive Goldstone mass as an adjustable input.

\section{Common tadpoles fix the Goldstone shift}
For any antisymmetric generator $T$ of a scalar symmetry and invariant
function $V$, differentiating $\nabla V\cdot TX=0$ gives
\begin{equation}
 V''(X)TX=T\nabla V(X).
 \label{eq:ward}
\end{equation}
This applies separately to the locally separated hard spectral
functional $V_{1,h}$ and to its invariant finite counterterm. On the
radial background, $\nabla V_{1,h}=RG^{-1}t_h$. The common tadpole
condition is $t_h+\partial_r V_{\rm CT}=0$. Consequently
\begin{equation}
 (V_{1,h}''+V_{\rm CT}'')TX=0.
 \label{eq:goldmass}
\end{equation}
Collecting the 45 gauge generators and the PQ generator into the
redundant orbit matrix $B=(T_aX)$, Eq.~\eqref{eq:ward} determines the
whole Goldstone block as
\begin{equation}
 B^T V_{1,h}'' B=B^T(T_a\nabla V_{1,h})_a.
\end{equation}
Its normalized 34-dimensional hard block has eigenvalues between
$-2.7348720\omega^2$ and $-0.0342377\omega^2$. The common counterterm
cancels it; the matrix residual is $7.8\times10^{-15}\omega^2$.
The actual tree cubic vertices independently obey the differentiated
identity
\begin{equation}
 Q^T T_r(T_aX)=Q^TT_a\mathsf H R_r,
 \label{eq:cubicward}
\end{equation}
where $Q$ spans the full tree kernel. This follows by differentiating
Eq.~\eqref{eq:ward} for $V_0$ and using $Q^T\mathsf H=0$.
All three radial derivatives are checked directly.

In a reorganized propagator the Goldstone shift must include its
counterterm in this convention. Resumming the bare hard shift alone,
or replacing it by a freely chosen positive regulator, would not
describe the same prescribed point. Equation~\eqref{eq:goldmass} is
a one-loop Ward result, not a completed two-loop resummation calculation.

\section{The actual soft infrared singularity}
\subsection{Gaussian determinant and the soft graph coefficient}
For real scalar fluctuations, Gaussian integration gives
\begin{equation}
 \Gamma_1^S=\tfrac12\Tr\log(-\partial_E^2+\mathsf H),\qquad
 V_1^S=\frac1{64\pi^2}\tr\left[
 \mathsf H^2\left(\log\frac{\mathsf H}{\mu_0^2}-\frac32\right)\right].
\end{equation}
Let $Q$ span the 38 soft modes and let $P_h$ span their massive
complement. Set
\begin{align}
 A_r&=Q^TT_rQ,\\
 \mathcal B_{rs}&=\frac{\tr(A_rA_s)}{32\pi^2},\label{eq:gram}\\
 \mathcal E_{rs}&=\frac1{32\pi^2}\left[
 \tr(Q^TU_{rs}Q)
 -2\sum_{a\in s,\,i\in h}
 \frac{(T_r)_{ai}(T_s)_{ia}}{m_i^2}\right].
\end{align}
The last term is essential: second-order eigenvalue perturbation
contains mixing with the hard complement. Freezing the soft projector
through second order would omit it.

Regulate the soft eigenvalue branches by adding a common mass squared
$\epsilon>0$, without changing the hard eigenvalues. Their tadpole
and Hessian at the reference point are
\begin{align}
 t^s_r(\epsilon)&=\frac{\epsilon}{32\pi^2}
 \left(\log\frac{\epsilon}{\mu_0^2}-1\right)\tr A_r
 \longrightarrow0,\label{eq:softtad}\\
 H^s_{rs}(\epsilon)&=\mathcal B_{rs}\log\frac{\epsilon}{\mu_0^2}
 +\epsilon\left(\log\frac{\epsilon}{\mu_0^2}-1\right)\mathcal E_{rs}.
 \label{eq:softH}
\end{align}
These follow from $f'(\epsilon)=2\epsilon(\log(\epsilon/\mu_0^2)-1)$,
$f''(\epsilon)=2\log(\epsilon/\mu_0^2)$ and the reduced second
eigenvalue derivative. Direct spectral finite differences independently
verify Eq.~\eqref{eq:softH}.

For a real radial vector $v$,
\begin{equation}
 v^T\mathcal Bv=\frac1{32\pi^2}
 \left\|\sum_r v_rA_r\right\|_F^2\geq0.
 \label{eq:psd}
\end{equation}
Thus the leading infrared logarithm is negative semidefinite as
$\epsilon\to0$. It is not a finite positive repair of a negative hard
curvature. In the declared radial coordinates, the actual matrix is
\begin{equation}
 \frac{\mathcal B}{\omega^2}\simeq
 \begin{pmatrix}
 0.60561772&0.00943837&0.00589905\\
 0.00943837&0.00048543&0.00020861\\
 0.00589905&0.00020861&0.00058130
 \end{pmatrix}.
\end{equation}
The smallness of its projection on one negative direction does not
mean that the whole matrix is small.

For the normalized negative witness $q$ retained from P-REN1,
$q^TGq=1$, its decomposition is
\begin{center}
\begin{tabular}{lr}
\toprule
Soft sector & $q^T\mathcal Bq/\omega^2$\\
\midrule
Gauge Goldstones & $3.1096973\times10^{-4}$\\
Physical PQ & $1.8983181\times10^{-6}$\\
Light Higgs & $2.6290000\times10^{-7}$\\
Total & $3.1313095\times10^{-4}$\\
\bottomrule
\end{tabular}
\end{center}
Mixed sector contributions vanish for these radial vertices, as
verified with the actual projectors. At $\epsilon/\omega^2=10^{-4}$
and $10^{-10}$, respectively, $q^TH^sq/\omega^2$ is approximately
$-0.00255$ and $-0.00685$. This is regulator dependence, not a pair
of physical mass predictions. The zero-momentum full potential Hessian
cannot be promoted to a pole mass by merely restoring the missing modes.

\section{Infrared cancellation at nonzero momentum}
\subsection{Complete one-loop scalar term on the radial slice}
Define the Euclidean finite bubble with $s=p_E^2>0$:
\begin{equation}
 L(a,b;s)=\int_0^1\log\frac{(1-x)a+xb+x(1-x)s}{\mu_0^2}\,dx.
 \label{eq:bubble}
\end{equation}
It is the negative of the usual finite $B_0$ in this convention.
Differentiating the scalar functional determinant twice gives
\begin{equation}
 \Pi^S_{rs}(s)=\frac1{32\pi^2}\left[
 \sum_i(U_{rs})_{ii}\,m_i^2\left(\log\frac{m_i^2}{\mu_0^2}-1\right)
 +\sum_{ij}(T_r)_{ij}(T_s)_{ji}L(m_i^2,m_j^2;s)\right].
 \label{eq:scalarself}
\end{equation}
The first term is the seagull and the second the two-cubic-vertex
bubble; ordered real-field indices account for the symmetry factor.
The sums include all 328 scalars, including mixed hard/soft pairs.
At $s=0$, removing only the soft--soft bubble reproduces the old scalar
hard Fr\'echet Hessian to relative residual $1.4\times10^{-16}$.

For two massless lines,
\begin{equation}
 L(0,0;s)=\log\frac{s}{\mu_0^2}-2.
 \label{eq:masslessbubble}
\end{equation}
For one massive line $m^2$ and $z=s/m^2$,
\begin{equation}
 L(0,m^2;s)=\log\frac{m^2}{\mu_0^2}-2
                +(1+z^{-1})\log(1+z).
\end{equation}
Independent Feynman-parameter integrals check both finite constants.
The logarithmic regulator cancels explicitly:
\begin{align}
 H^s(\epsilon)&+\mathcal B
 \left[L(\epsilon,\epsilon;s)-\log\frac{\epsilon}{\mu_0^2}\right]
 \nonumber\\
 &\longrightarrow\mathcal B\left(\log\frac{s}{\mu_0^2}-2\right).
 \label{eq:cancel}
\end{align}
At fixed $s>0$ this is finite. It does not license taking a physical
mass from the $s=0$ potential curvature.

For at least one massive line,
\begin{equation}
 L(a,b;s)-L(a,b;0)
 =\int_0^1\log\left[1+\frac{x(1-x)s}{(1-x)a+xb}\right]dx\geq0.
\end{equation}
Combined with the vertex Gram form, the hard scalar momentum
difference $\Delta\Pi^S_h(s)$ is positive semidefinite. By contrast,
the massless contribution is negative semidefinite for
$0<s\leq e^2\mu_0^2$ in this MS prescription.

The displayed diagnostic kernel is
\begin{equation}
 K_{\rm diag}(s)=sG+H_{r,B}
 +\Delta\Pi^S_h(s)+\mathcal B[\log(s/\mu_0^2)-2],
 \label{eq:diag}
\end{equation}
where $H_{r,B}$ is the old bosonic hard-plus-counterterm radial block.
The vector contribution is still evaluated at zero momentum.
\begin{center}
\begin{tabular}{rrrrrrr}
\toprule
$s/\omega^2$ & $.001$ & $.01$ & $.05$ & $.1$ & $.5$ & $1$\\
\midrule
$\lambda_{\min}(K_{\rm diag},G)/\omega^2$
& $-.75843$ & $-.72498$ & $-.58421$ & $-.41369$ & $.59623$ & $1.09859$\\
\bottomrule
\end{tabular}
\end{center}
These are kinematic probes, not a parameter scan or physical pole
locations. Missing vector/Goldstone/ghost mixed diagrams, fermion
momentum maps and Nielsen consistency prevent interpreting a zero of
this partial kernel as a physical tachyon or a stable-particle pole.

\subsection{A separate quantitative failure of small-insertion control}
The finite hard scalar kinetic matrix is
\begin{equation}
 (\Delta Z_h^S)_{rs}=\frac1{32\pi^2}
 \sum_{(i,j)\notin(s,s)}(T_r)_{ij}(T_s)_{ji}
 \int_0^1\frac{x(1-x)\,dx}{(1-x)m_i^2+x m_j^2}.
 \label{eq:kinetic}
\end{equation}
It is positive semidefinite and its generalized eigenvalues relative
to $G$ are
\begin{equation}
 (0.00900848,\;0.27686405,\;2.61031989).
\end{equation}
Finite differences of the actual momentum kernel verify this derivative.
The soft--soft contribution is nonlocal and is not included in
Eq.~\eqref{eq:kinetic}. The largest hard correction already exceeds
the tree metric, so the naive small-insertion expansion fails this
necessary engineering check at the frozen point.

A finite canonical field transformation can absorb a positive metric,
but is not by itself a bound on omitted loop orders. Moreover, an
invertible congruence of a fixed curvature matrix preserves its inertia.
It cannot turn its negative direction positive simply by renormalizing
the fields. These statements are not an all-orders exclusion.

\section{An entire simple repair ray fails analytically}
Uniformly multiply every existing scalar coefficient by $z>0$, including
the dependent masses and the cubic coefficient, keeping $r_0,g_U,\mu_0$
fixed. Then $V_0\mapsto zV_0$ exactly: stationarity, the tree kernel and
tree inertia are unchanged, while $\mathsf H,T_r,U_{rs}$ each scale by $z$.
The common fixed-VEV scalar contribution therefore obeys
\begin{equation}
 H_h(z)=zH_0+z^2[S+\log z\,L_S]+V,
 \qquad \Delta Z_h^S(z)=z\Delta Z_h^S(1).
 \label{eq:ray}
\end{equation}
Here $S$ includes the scalar loop and its own share of the common
counterterm, $V$ the vector loop and its share, and $L_S$ the coefficient
of the scalar logarithm including its tadpole conversion. Direct
rescaled Fr\'echet calculations check Eq.~\eqref{eq:ray}; the following
bound is not inferred from the sampled points.

\begin{proposition}
For every $0<z\leq1$, the bosonic hard fixed-VEV radial block on this
ray has a negative direction.
\end{proposition}
\begin{proof}
In the sigma coordinate the actual coefficients, in units of
$\omega^2$, are
\begin{equation}
 \begin{aligned}
 a=(H_0)_{\sigma\sigma}=0.02880405,\quad
 &b=-S_{\sigma\sigma}=1.54130008,\\
 (L_S)_{\sigma\sigma}=7.61695881,\quad
 &c=V_{\sigma\sigma}=-0.00402215.
 \end{aligned}
\end{equation}
Since $\log z\leq0$, one has
\begin{equation}
 (H_h(z))_{\sigma\sigma}/\omega^2
 \leq az-bz^2+c\leq c+\frac{a^2}{4b}
 =-0.00388757<0.
\end{equation}
The normalized radial direction is $e_\sigma/\sqrt2$, and remains
negative for the entire interval.
\end{proof}
The theorem concerns this bosonic ray only. It does not exclude
nonuniform scalar choices, different radial hierarchies, or fermionic
corrections at new scalar points. It shows why merely reducing all
scalar couplings is not enough in the fixed gauge/scale convention.

\section{Bounded reselection without fitting observations}
After the original control failure, four coarse engineering points
were tested with $(z,\sigma/\omega)=(.03,.3),(.03,.5),(.1,.3),(.1,.5)$
and $w/\omega=1$, $v_s/\omega=.25$. All existing scalar coefficients
are first scaled by $z$. The dependent masses are then re-solved from
\begin{equation}
 A(r)\,\delta p=-\partial_r V_0(r;zp),\qquad
 A(r)=-\diag(12w/5,\sigma,v_s).
\end{equation}
The tree $\xi_{02}$ condition is reimposed at the first physical
doublet crossing. Every point has 34 symmetry zero modes plus precisely
one complex light doublet, checked with the actual SM Casimirs.

The polynomial identity
\begin{equation}
 \mathsf H(r_0+\Delta\sigma e_\sigma;p)
 =\mathsf H_0+\Delta\sigma T_\sigma
   +\tfrac12(\Delta\sigma)^2 U_{\sigma\sigma}
\end{equation}
and the exact mass-parameter operator give the new full ambient tree
Hessian without a fresh large tensor differentiation. New tadpoles,
common counterterms and radial hard/kinetic matrices are then evaluated
from each point's own spectrum, not borrowed from the old point.

We demand radial hard positivity and test two dimensionless control
quantities, with a declared engineering margin $0.3$:
\begin{equation}
 \rho_H=\max|\operatorname{spec}(H_h-H_0,H_0)|,
 \qquad \rho_Z=\max\operatorname{spec}(\Delta Z_h^S,G).
\end{equation}
This margin is not a theorem of perturbative convergence; importantly,
the rejected positive point even fails the weaker $\rho_H<1$ test.
\begin{center}
\begin{tabular}{rrrrr}
\toprule
$z$ & $\sigma/\omega$ & $\lambda_{\min}(H_h,G)/\omega^2$ & $\rho_H$ & $\rho_Z$\\
\midrule
$.03$ & $.3$ & $-.040838$ & $18.8648$ & $.09010$\\
$.03$ & $.5$ & $-.044064$ & $7.9561$ & $.09296$\\
$.1$ & $.3$ & $-.198595$ & $27.0000$ & $.30032$\\
$.1$ & $.5$ & $.009714$ & $1.4323$ & $.30988$\\
\bottomrule
\end{tabular}
\end{center}
Thus none is accepted. The last point illustrates why positivity alone
cannot certify a controlled background. This is not a no-go theorem
for all less-hierarchical points. Full parameter cards, finite tadpole
shifts, tree classifications and rejection evidence are retained in
the reselection JSON. No measured mass or flavor datum enters.
Changing the hierarchy invalidates reuse of the old two-loop
unification/threshold solution; $g_U$ was held only as a reference
coupling for these engineering probes.

\section{Status and next mathematical tasks}
The IR verifier passes 65/65 checks and the bounded reselection verifier
passes 36/36. Passing includes correctly rejecting physical promotion.
Both reuse ten cached Hessians and require no new Hessian evaluations.
The new artifacts are \path{code/verify_p54_goldstone_ir.py} and
\path{code/verify_p54_scalar_reselection.py}, with matching JSON/Markdown
reports under \path{output/}; sources and hashes are recorded.

The useful next options are now more specific:
\begin{enumerate}
\item Decompose the actual scalar-invariant contributions to
$\rho_H$ and $\rho_Z$, especially the independent $126$ quartics,
before nonuniform parameter adjustment. Bare couplings smaller than
one and positive selected masses are not sufficient control criteria.
Do not fine-tune against the displayed Euclidean subset.
\item Complete the same-prescription gauge/Goldstone/ghost and fermion
momentum maps with Nielsen consistency if deciding the physical fate
of the original point. A true pole/background claim cannot follow
from a scalar-only repair or the partial kernel in Eq.~\eqref{eq:diag}.
\item After finding a controlled common scalar point, recompute its
full nonradial response, one-loop light state, upper PS saddle and
threshold/running system. Complete Wilson/box, lower finite matching
and finite seesaw remain independent unfinished calculations. None
is promoted to a physical fit by this audit.
\end{enumerate}

\begin{thebibliography}{9}
\bibitem{bg} J. Braathen and M. D. Goodsell,
Avoiding the Goldstone Boson Catastrophe in general renormalisable field
theories at two loops, especially the distinction between tadpole
resummation and momentum-regulated masses,
\href{https://arxiv.org/html/1609.06977v2}{arXiv:1609.06977v2}.
The two-loop mass results in that work specialize to the gaugeless limit;
they are not a completed P54 gauge-sector calculation.
\bibitem{ek} J. R. Espinosa and T. Konstandin,
Resummation of Goldstone Infrared Divergences: A Proof to All Orders,
\href{https://arxiv.org/abs/1712.08068}{arXiv:1712.08068}.
\end{thebibliography}
\end{document}
